\scanpage{27}%
\mspage{12}%
\margin{\handwritten{$\begin{aligned}j&=0 &&\Longleftrightarrow c_4=0\\j&=12^3 &&\Longleftrightarrow c_6=0\end{aligned}$}}%
\[
j=c_4^3/\Delta=1728+c_6^2/\Delta \, . \tag*{(6.2)}
\]

\uline{Proposition} 6.3. \quad \uline{Soient} $E$ \uline{et} $E'$ \uline{deux courbes de genre un sur} $S$, \uline{sans fibre additive} .

(I) $\text{\uline{Isom}}_{\handwritten{S}}(E,E')$ \uline{est représentable, fini et} \tikz[baseline=(nrone.base)]{\node[inner sep=0pt] (nrone) {\uline{non ramifié}};\handwritten{\draw[thin,gray] (nrone.north east) .. controls ([xshift=-1em,yshift=0.8em]nrone.north east) and ([xshift=-0.6em,yshift=0.7em]nrone.north west) .. ([xshift=-0.2em]nrone.west) .. controls ([yshift=-0.85em]nrone.south west) and ([xshift=0.5em,yshift=-0.85em]nrone.south east) .. (nrone.north east);}\node[anchor=south west,inner sep=0pt] at ([xshift=-1.8em,yshift=0.48em]nrone.north east) {\inserted{\handwritten{net}}};} \uline{sur} $S$\tikz[baseline=(stop.base)]{\node[inner sep=0pt] (stop) {.};\handwritten{\draw[thin] ([yshift=-0.08em]stop.south) .. controls ([xshift=0.04em,yshift=-0.19em]stop.south) .. ([xshift=-0.15em,yshift=-0.35em]stop.south);\draw[thin,gray,rotate around={-18:(stop.center)}] (stop.center) ellipse (0.56em and 1.34em);}}

(II) \uline{si} $j_E=j_F$, \uline{la projection de} $\text{\uline{I\mbox{$\mathbb{\text{\makebox[0pt][l]{\kern.45pt$s$}$s$}}$}om}}_{\handwritten{S}}(E,E')$ \uline{sur} $S$ \uline{est surjective}\tikz[baseline=(stop.base)]{\node[inner sep=0pt] (stop) {.};\handwritten{\draw[thin] ([yshift=-0.08em]stop.south) .. controls ([xshift=0.04em,yshift=-0.19em]stop.south) .. ([xshift=-0.15em,yshift=-0.35em]stop.south);\draw[thin,gray,rotate around={-18:(stop.center)}] (stop.center) ellipse (0.56em and 1.34em);}}

(III) \uline{si} $j_E=j_F$ \uline{ne prend pas les valeurs} $0$ \uline{et} $\underset{\inserted{\handwritten{\substack{\Vert\\12^3}}}}{\tikz[baseline=(val.base)]{\node[inner sep=0pt] (val) {$1728$};\handwritten{\draw[thin] (val.south west) .. controls ([yshift=-0.15em]val.south west) and ([yshift=-0.15em]val.south east) .. (val.south east);}}}$, $\text{\uline{Isom}}_{\handwritten{S}}(E,E')$ \uline{est un revêtement étale de rang} $2$ \uline{de} $S$ .

Donnons-nous $E$ et $E'$ par des équations (1.3) et $(1.3)'$. Se donner un isomorphisme entre $E$ et $E'$ revient à se donner $u$, $r$, $s$ et $t$, vérifiant les équations (1.6), donc aussi (1.7) et (1.8), $u$ devant de plus être inversible. Par hypothèse, localement sur $S$, $c_4$ ou $\Delta$ est inversible (6.1), et de même $c'_4$ ou $\Delta'$ l'est. Localement sur $S$, les équations (1.8) sont donc des équations de dépendance intégrale pour $u$ et $u^{-1}$, les équations $(1.7)_6$ ou $(1.7)_8$ des équations de dépendance intégrale pour $r$ et enfin $(1.6)_6$ (resp. $(1.6)_2$) est une équation de dépendance intégrale pour $t$ (resp. $s$). Ceci prouve que $\text{\uline{Isom}}(E,E')$ est représentable par un schéma affine fini sur $S$ ; il est même \tikz[baseline=(nrtwo.base)]{\node[inner sep=0pt] (nrtwo) {non ramifié};\handwritten{\draw[thin,gray] (nrtwo.north east) .. controls ([xshift=-1em,yshift=0.8em]nrtwo.north east) and ([xshift=-0.6em,yshift=0.7em]nrtwo.north west) .. ([xshift=-0.2em]nrtwo.west) .. controls ([yshift=-0.85em]nrtwo.south west) and ([xshift=0.5em,yshift=-0.85em]nrtwo.south east) .. (nrtwo.north east);}\node[anchor=south west,inner sep=0pt] at ([xshift=-1.8em,yshift=0.48em]nrtwo.north east) {\inserted{\handwritten{net}}};} car les courbes considérées n'ont pas d'automorphismes infinitésimaux (respectant l'origine).

Lorsque $2$ est inversible sur $S$, l'assertion (III) résulte de l'énoncé plus précis.

\uline{Proposition} 6.4. \quad \uline{Soient} $E$ \uline{et} $E'$ \uline{deux courbes de genre un sur} $S$ ,%
% archive sheet number: 26
